\documentclass[12pt]{article}
\usepackage[margin=0.72in]{geometry}
\usepackage{lmodern}
\usepackage{microtype}
\usepackage{multicol}
\usepackage{fancyhdr}
\usepackage{titlesec}
\usepackage{enumitem}
\usepackage{xcolor}
\usepackage[hidelinks]{hyperref}

\setlength{\columnsep}{0.30in}
\setlength{\parindent}{1.05em}
\setlength{\parskip}{0.24em}
\setlength{\headheight}{15pt}
\titleformat{\section}{\large\bfseries\scshape}{}{0pt}{}
\titleformat{\subsection}{\normalsize\bfseries}{}{0pt}{}
\pagestyle{fancy}
\fancyhf{}
\fancyhead[L]{Homer's Sirens}
\fancyhead[R]{Reading Handout}
\fancyfoot[C]{\thepage}
\renewcommand{\headrulewidth}{0.4pt}

\newcommand{\focusbox}[1]{%
\vspace{0.4em}\noindent\colorbox{gray!12}{\begin{minipage}{0.96\linewidth}#1\end{minipage}}\vspace{0.5em}}

\begin{document}
\begin{center}
{\LARGE\bfseries Do We Love Shadows?}\par\vspace{0.25em}
{\large\scshape Reading 2: Homer and the Sirens}\par\vspace{0.45em}
\rule{0.82\textwidth}{0.4pt}\par\vspace{0.7em}
\end{center}

\section*{Vocabulary Preview}
\begin{multicols}{2}
\begin{itemize}[leftmargin=*, itemsep=0.18em]
\item \textbf{Siren}: A dangerous singing creature who lures sailors toward destruction.
\item \textbf{Enchant}: To charm, delight, or capture the mind powerfully.
\item \textbf{Renowned}: Famous and honored.
\item \textbf{Achaeans}: A name Homer uses for the Greeks.
\item \textbf{Prophecy}: A warning or message about what will happen.
\item \textbf{Companion}: A fellow traveler or friend.
\item \textbf{Mast}: The tall upright pole that holds a ship's sail.
\item \textbf{Lash}: To tie tightly with rope.
\item \textbf{Restrain}: To hold back or prevent someone from acting.
\item \textbf{Temptation}: An attractive invitation toward something dangerous or wrong.
\item \textbf{Curiosity}: The desire to know, see, or experience something.
\item \textbf{Self-command}: The ability to govern one's own desires and actions.
\end{itemize}
\end{multicols}

\section*{Introduction}
\textit{The Odyssey} is an ancient Greek epic poem about Odysseus's long journey home from the Trojan War. In Book XII, Odysseus tells how the goddess Circe warned him about several dangers he would meet at sea. The first danger is the Sirens, whose song is so beautiful and persuasive that sailors forget home, family, duty, and life itself.

The Sirens are not a simple test of pleasure. They promise knowledge. They claim that anyone who listens will become wiser. This makes their song especially dangerous for a unit about truth and shadows: sometimes a shadow does not look ugly or foolish. Sometimes it looks beautiful, intelligent, and worth hearing. As you read, ask: \textbf{Why would Odysseus want to hear a song he knows may destroy him?}

\focusbox{\textbf{Source note:} Adapted and modernized from Homer, \textit{The Odyssey}, Book XII, in Samuel Butler's public-domain translation. This packet uses the focused Sirens episode: Circe's warning, Odysseus's instructions to his crew, and the passing of the Sirens.}

\begin{multicols}{2}
\section*{Reading}

\subsection*{1. Circe's Warning}
When morning came, Circe spoke to Odysseus about the dangers ahead.

"First," she said, "you will come to the Sirens. They enchant every man who comes near them. If a sailor draws too close and hears their song, he will never again be welcomed home by his wife and children. The Sirens sit in a meadow, singing sweetly, but around them lie the bones of men whose flesh has wasted away.

"Therefore you must sail past them. Soften beeswax and stop the ears of your men so that none of them can hear. But if you wish to listen, you may. Have your men bind you upright to the mast, with ropes fastened tightly around you. Then you may hear the Sirens' voices as the ship passes.

"But if you beg and command your men to release you, they must bind you with still more ropes."

\subsection*{2. Odysseus Warns His Crew}
After leaving Circe's island, Odysseus spoke to his companions.

"Friends," he said, "it is not right that only one or two of us should know the warnings Circe gave me. I will tell you what she said, so that whether we live or die, we may know what danger is coming.

"First, she said we must keep clear of the Sirens and their song. She ordered all of you to stop your ears with wax. But she said that I alone may hear them, if you bind me tightly to the mast. If I beg you or order you to set me free, you must only tie me more firmly."

While he spoke, the ship moved swiftly over the water. Soon they were near the Sirens' island.

\subsection*{3. The Crew Is Made Deaf}
The wind suddenly dropped, and the sea grew calm. Some god had stilled the waves. Odysseus's men took down the sail and laid it in the ship. Then they sat at the oars and struck the water with their blades.

Odysseus took a large round of wax, cut it into pieces with his sword, and kneaded the pieces in his strong hands until the sun's heat softened them. Then he pressed the wax into the ears of each man.

After that, the men bound Odysseus by hands and feet, standing upright against the mast. They lashed the ropes to the mast itself. Then they sat down again and rowed.

\subsection*{4. The Sirens Sing}
When the ship came within hearing distance, the Sirens saw it driving past and began their clear song:

"Come this way, renowned Odysseus, great honor of the Achaeans. Stop your ship so that you may hear our voices. No one has ever sailed past us without stopping to listen to the sweetness of our song. Whoever listens goes on his way delighted and wiser than before.

"We know everything the Greeks and Trojans suffered at Troy by the will of the gods. We know all that happens on the fruitful earth."

Their voices poured over the water, and Odysseus longed to hear more. He frowned and signaled to his men with his eyebrows, ordering them to release him.

But the men bent harder to their rowing. Two of them, Perimedes and Eurylochus, rose from their seats and tied Odysseus with more ropes, binding him still more tightly to the mast.

\subsection*{5. Passing the Danger}
At last the ship had passed beyond the Sirens. Their song grew faint behind them. When the men could no longer hear it, they took the wax from their ears and untied Odysseus from the mast.

The danger was over, but only because Odysseus had known two things at once: he wanted to hear the song, and he could not trust himself to master that desire when the song began.

\section*{Reading Focus}
\begin{enumerate}[leftmargin=*, itemsep=0.2em]
\item Underline the sentence in Circe's warning that shows the Sirens are deadly.
\item Circle the part of the Sirens' song where they promise knowledge or wisdom.
\item Put a star beside the moment when Odysseus shows that desire can overpower judgment.
\item Write "self-command" in the margin beside the action that keeps Odysseus safe.
\item Box one detail showing that Odysseus needs help from other people, not just his own willpower.
\end{enumerate}

\section*{Window Note}
Create a three-pane window note. In the first pane, draw what the Sirens promise. In the second, draw what their song hides. In the third, draw Odysseus bound to the mast while his men row past. Under each pane, write one sentence explaining what the scene shows about desire, truth, illusion, or self-command.

\section*{Claim Building}
Answer in one clear paragraph: \textbf{Why does Odysseus want to hear the Sirens if he knows they are dangerous?} State a claim, use one specific moment from the episode, explain your evidence, and connect your answer to the unit question.

\section*{Handout Summary}
In 3--5 sentences, summarize the Sirens episode. Include Circe's warning, the Sirens' promise, Odysseus's desire to be released, and the reason his crew must bind him tighter. End with one sentence explaining how Homer helps us understand why people may love shadows.

\section*{Comparison Seed}
Keep this sentence for later readings: \textbf{For Homer, people are drawn toward shadows because} \underline{\hspace{1.8in}}. Compare this with Plato: are the Sirens like the cave's shadows, or are they a different kind of danger?

\end{multicols}
\end{document}
